%!TeX root=../persuasiontop.tex
\chapter[Chapter \thechapter]{}

\begin{figure}[t!]
\centering
\includegraphics[width=\linewidth]{nouveau06}
\end{figure}

\lettrine[lines=4,lraise=0.3]{A}{nne} had not wanted this visit to Uppercross, to learn that a removal from one set of people to another, though at a distance of only three miles, will often include a total change of conversation, opinion, and idea. She had never been staying there before, without being struck by it, or without wishing that other Elliots could have her advantage in seeing how unknown, or unconsidered there, were the affairs which at Kellynch Hall were treated as of such general publicity and pervading interest; yet, with all this experience, she believed she must now submit to feel that another lesson, in the art of knowing our own nothingness beyond our own circle, was become necessary for her; for certainly, coming as she did, with a heart full of the subject which had been completely occupying both houses in Kellynch for many weeks, she had expected rather more curiosity and sympathy than she found in the separate but very similar remark of Mr and Mrs Musgrove: »So, Miss Anne, Sir Walter and your sister are gone; and what part of Bath do you think they will settle in?« and this, without much waiting for an answer; or in the young ladies' addition of, »I hope \textit{we} shall be in Bath in the winter; but remember, papa, if we do go, we must be in a good situation: none of your Queen Squares for us!« or in the anxious supplement from Mary, of—»Upon my word, I shall be pretty well off, when you are all gone away to be happy at Bath!«

She could only resolve to avoid such self-delusion in future, and think with heightened gratitude of the extraordinary blessing of having one such truly sympathising friend as Lady Russell.

The Mr Musgroves had their own game to guard, and to destroy, their own horses, dogs, and newspapers to engage them, and the females were fully occupied in all the other common subjects of housekeeping, neighbours, dress, dancing, and music. She acknowledged it to be very fitting, that every little social commonwealth should dictate its own matters of discourse; and hoped, ere long, to become a not unworthy member of the one she was now transplanted into. With the prospect of spending at least two months at Uppercross, it was highly incumbent on her to clothe her imagination, her memory, and all her ideas in as much of Uppercross as possible.

She had no dread of these two months. Mary was not so repulsive and unsisterly as Elizabeth, nor so inaccessible to all influence of hers; neither was there anything among the other component parts of the cottage inimical to comfort. She was always on friendly terms with her brother-in-law; and in the children, who loved her nearly as well, and respected her a great deal more than their mother, she had an object of interest, amusement, and wholesome exertion.

Charles Musgrove was civil and agreeable; in sense and temper he was undoubtedly superior to his wife, but not of powers, or conversation, or grace, to make the past, as they were connected together, at all a dangerous contemplation; though, at the same time, Anne could believe, with Lady Russell, that a more equal match might have greatly improved him; and that a woman of real understanding might have given more consequence to his character, and more usefulness, rationality, and elegance to his habits and pursuits. As it was, he did nothing with much zeal, but sport; and his time was otherwise trifled away, without benefit from books or anything else. He had very good spirits, which never seemed much affected by his wife's occasional lowness, bore with her unreasonableness sometimes to Anne's admiration, and upon the whole, though there was very often a little disagreement (in which she had sometimes more share than she wished, being appealed to by both parties), they might pass for a happy couple. They were always perfectly agreed in the want of more money, and a strong inclination for a handsome present from his father; but here, as on most topics, he had the superiority, for while Mary thought it a great shame that such a present was not made, he always contended for his father's having many other uses for his money, and a right to spend it as he liked.

As to the management of their children, his theory was much better than his wife's, and his practice not so bad. »I could manage them very well, if it were not for Mary's interference,« was what Anne often heard him say, and had a good deal of faith in; but when listening in turn to Mary's reproach of »Charles spoils the children so that I cannot get them into any order,« she never had the smallest temptation to say, »Very true.«

One of the least agreeable circumstances of her residence there was her being treated with too much confidence by all parties, and being too much in the secret of the complaints of each house. Known to have some influence with her sister, she was continually requested, or at least receiving hints to exert it, beyond what was practicable. »I wish you could persuade Mary not to be always fancying herself ill,« was Charles's language; and, in an unhappy mood, thus spoke Mary: »I do believe if Charles were to see me dying, he would not think there was anything the matter with me. I am sure, Anne, if you would, you might persuade him that I really am very ill—a great deal worse than I ever own.«

Mary's declaration was, »I hate sending the children to the Great House, though their grandmamma is always wanting to see them, for she humours and indulges them to such a degree, and gives them so much trash and sweet things, that they are sure to come back sick and cross for the rest of the day.« And Mrs Musgrove took the first opportunity of being alone with Anne, to say, »Oh! Miss Anne, I cannot help wishing Mrs Charles had a little of your method with those children. They are quite different creatures with you! But to be sure, in general they are so spoilt! It is a pity you cannot put your sister in the way of managing them. They are as fine healthy children as ever were seen, poor little dears! without partiality; but Mrs Charles knows no more how they should be treated—! Bless me! how troublesome they are sometimes. I assure you, Miss Anne, it prevents my wishing to see them at our house so often as I otherwise should. I believe Mrs Charles is not quite pleased with my not inviting them oftener; but you know it is very bad to have children with one that one is obligated to be checking every moment; »don't do this,« and »don't do that;« or that one can only keep in tolerable order by more cake than is good for them.«

She had this communication, moreover, from Mary. »Mrs Musgrove thinks all her servants so steady, that it would be high treason to call it in question; but I am sure, without exaggeration, that her upper house-maid and laundry-maid, instead of being in their business, are gadding about the village, all day long. I meet them wherever I go; and I declare, I never go twice into my nursery without seeing something of them. If Jemima were not the trustiest, steadiest creature in the world, it would be enough to spoil her; for she tells me, they are always tempting her to take a walk with them.« And on Mrs Musgrove's side, it was, »I make a rule of never interfering in any of my daughter-in-law's concerns, for I know it would not do; but I shall tell \textit{you}, Miss Anne, because you may be able to set things to rights, that I have no very good opinion of Mrs Charles's nursery-maid: I hear strange stories of her; she is always upon the gad; and from my own knowledge, I can declare, she is such a fine-dressing lady, that she is enough to ruin any servants she comes near. Mrs Charles quite swears by her, I know; but I just give you this hint, that you may be upon the watch; because, if you see anything amiss, you need not be afraid of mentioning it.«

Again, it was Mary's complaint, that Mrs Musgrove was very apt not to give her the precedence that was her due, when they dined at the Great House with other families; and she did not see any reason why she was to be considered so much at home as to lose her place. And one day when Anne was walking with only the Musgroves, one of them after talking of rank, people of rank, and jealousy of rank, said, »I have no scruple of observing to \textit{you}, how nonsensical some persons are about their place, because all the world knows how easy and indifferent you are about it; but I wish anybody could give Mary a hint that it would be a great deal better if she were not so very tenacious, especially if she would not be always putting herself forward to take place of mamma. Nobody doubts her right to have precedence of mamma, but it would be more becoming in her not to be always insisting on it. It is not that mamma cares about it the least in the world, but I know it is taken notice of by many persons.«

How was Anne to set all these matters to rights? She could do little more than listen patiently, soften every grievance, and excuse each to the other; give them all hints of the forbearance necessary between such near neighbours, and make those hints broadest which were meant for her sister's benefit.

In all other respects, her visit began and proceeded very well. Her own spirits improved by change of place and subject, by being removed three miles from Kellynch; Mary's ailments lessened by having a constant companion, and their daily intercourse with the other family, since there was neither superior affection, confidence, nor employment in the cottage, to be interrupted by it, was rather an advantage. It was certainly carried nearly as far as possible, for they met every morning, and hardly ever spent an evening asunder; but she believed they should not have done so well without the sight of Mr and Mrs Musgrove's respectable forms in the usual places, or without the talking, laughing, and singing of their daughters.

She played a great deal better than either of the Miss Musgroves, but having no voice, no knowledge of the harp, and no fond parents, to sit by and fancy themselves delighted, her performance was little thought of, only out of civility, or to refresh the others, as she was well aware. She knew that when she played she was giving pleasure only to herself; but this was no new sensation. Excepting one short period of her life, she had never, since the age of fourteen, never since the loss of her dear mother, known the happiness of being listened to, or encouraged by any just appreciation or real taste. In music she had been always used to feel alone in the world; and Mr and Mrs Musgrove's fond partiality for their own daughters' performance, and total indifference to any other person's, gave her much more pleasure for their sakes, than mortification for her own.

The party at the Great House was sometimes increased by other company. The neighbourhood was not large, but the Musgroves were visited by everybody, and had more dinner-parties, and more callers, more visitors by invitation and by chance, than any other family. They were more completely popular.

The girls were wild for dancing; and the evenings ended, occasionally, in an unpremeditated little ball. There was a family of cousins within a walk of Uppercross, in less affluent circumstances, who depended on the Musgroves for all their pleasures: they would come at any time, and help play at anything, or dance anywhere; and Anne, very much preferring the office of musician to a more active post, played country dances to them by the hour together; a kindness which always recommended her musical powers to the notice of Mr and Mrs Musgrove more than anything else, and often drew this compliment;—»Well done, Miss Anne! very well done indeed! Lord bless me! how those little fingers of yours fly about!«

So passed the first three weeks. Michaelmas came; and now Anne's heart must be in Kellynch again. A beloved home made over to others; all the precious rooms and furniture, groves, and prospects, beginning to own other eyes and other limbs! She could not think of much else on the 29th of September; and she had this sympathetic touch in the evening from Mary, who, on having occasion to note down the day of the month, exclaimed, »Dear me, is not this the day the Crofts were to come to Kellynch? I am glad I did not think of it before. How low it makes me!«

The Crofts took possession with true naval alertness, and were to be visited. Mary deplored the necessity for herself. »Nobody knew how much she should suffer. She should put it off as long as she could;« but was not easy till she had talked Charles into driving her over on an early day, and was in a very animated, comfortable state of imaginary agitation, when she came back. Anne had very sincerely rejoiced in there being no means of her going. She wished, however, to see the Crofts, and was glad to be within when the visit was returned. They came: the master of the house was not at home, but the two sisters were together; and as it chanced that Mrs Croft fell to the share of Anne, while the Admiral sat by Mary, and made himself very agreeable by his good-humoured notice of her little boys, she was well able to watch for a likeness, and if it failed her in the features, to catch it in the voice, or in the turn of sentiment and expression.

Mrs Croft, though neither tall nor fat, had a squareness, uprightness, and vigour of form, which gave importance to her person. She had bright dark eyes, good teeth, and altogether an agreeable face; though her reddened and weather-beaten complexion, the consequence of her having been almost as much at sea as her husband, made her seem to have lived some years longer in the world than her real eight-and-thirty. Her manners were open, easy, and decided, like one who had no distrust of herself, and no doubts of what to do; without any approach to coarseness, however, or any want of good humour. Anne gave her credit, indeed, for feelings of great consideration towards herself, in all that related to Kellynch, and it pleased her: especially, as she had satisfied herself in the very first half minute, in the instant even of introduction, that there was not the smallest symptom of any knowledge or suspicion on Mrs Croft's side, to give a bias of any sort. She was quite easy on that head, and consequently full of strength and courage, till for a moment electrified by Mrs Croft's suddenly saying,—

»It was you, and not your sister, I find, that my brother had the pleasure of being acquainted with, when he was in this country.«

Anne hoped she had outlived the age of blushing; but the age of emotion she certainly had not.

»Perhaps you may not have heard that he is married?« added Mrs Croft.

She could now answer as she ought; and was happy to feel, when Mrs Croft's next words explained it to be Mr Wentworth of whom she spoke, that she had said nothing which might not do for either brother. She immediately felt how reasonable it was, that Mrs Croft should be thinking and speaking of Edward, and not of Frederick; and with shame at her own forgetfulness applied herself to the knowledge of their former neighbour's present state with proper interest.

The rest was all tranquillity; till, just as they were moving, she heard the Admiral say to Mary—

»We are expecting a brother of Mrs Croft's here soon; I dare say you know him by name.«

He was cut short by the eager attacks of the little boys, clinging to him like an old friend, and declaring he should not go; and being too much engrossed by proposals of carrying them away in his coat pockets, \&c., to have another moment for finishing or recollecting what he had begun, Anne was left to persuade herself, as well as she could, that the same brother must still be in question. She could not, however, reach such a degree of certainty, as not to be anxious to hear whether anything had been said on the subject at the other house, where the Crofts had previously been calling.

The folks of the Great House were to spend the evening of this day at the Cottage; and it being now too late in the year for such visits to be made on foot, the coach was beginning to be listened for, when the youngest Miss Musgrove walked in. That she was coming to apologize, and that they should have to spend the evening by themselves, was the first black idea; and Mary was quite ready to be affronted, when Louisa made all right by saying, that she only came on foot, to leave more room for the harp, which was bringing in the carriage.

»And I will tell you our reason,« she added, »and all about it. I am come on to give you notice, that papa and mamma are out of spirits this evening, especially mamma; she is thinking so much of poor Richard! And we agreed it would be best to have the harp, for it seems to amuse her more than the piano-forte. I will tell you why she is out of spirits. When the Crofts called this morning, (they called here afterwards, did not they?), they happened to say, that her brother, Captain Wentworth, is just returned to England, or paid off, or something, and is coming to see them almost directly; and most unluckily it came into mamma's head, when they were gone, that Wentworth, or something very like it, was the name of poor Richard's captain at one time; I do not know when or where, but a great while before he died, poor fellow! And upon looking over his letters and things, she found it was so, and is perfectly sure that this must be the very man, and her head is quite full of it, and of poor Richard! So we must be as merry as we can, that she may not be dwelling upon such gloomy things.«

The real circumstances of this pathetic piece of family history were, that the Musgroves had had the ill fortune of a very troublesome, hopeless son; and the good fortune to lose him before he reached his twentieth year; that he had been sent to sea because he was stupid and unmanageable on shore; that he had been very little cared for at any time by his family, though quite as much as he deserved; seldom heard of, and scarcely at all regretted, when the intelligence of his death abroad had worked its way to Uppercross, two years before.

He had, in fact, though his sisters were now doing all they could for him, by calling him »poor Richard,« been nothing better than a thick-headed, unfeeling, unprofitable Dick Musgrove, who had never done anything to entitle himself to more than the abbreviation of his name, living or dead.

He had been several years at sea, and had, in the course of those removals to which all midshipmen are liable, and especially such midshipmen as every captain wishes to get rid of, been six months on board Captain Frederick Wentworth's frigate, the Laconia; and from the Laconia he had, under the influence of his captain, written the only two letters which his father and mother had ever received from him during the whole of his absence; that is to say, the only two disinterested letters; all the rest had been mere applications for money.

In each letter he had spoken well of his captain; but yet, so little were they in the habit of attending to such matters, so unobservant and incurious were they as to the names of men or ships, that it had made scarcely any impression at the time; and that Mrs Musgrove should have been suddenly struck, this very day, with a recollection of the name of Wentworth, as connected with her son, seemed one of those extraordinary bursts of mind which do sometimes occur.

She had gone to her letters, and found it all as she supposed; and the re-perusal of these letters, after so long an interval, her poor son gone for ever, and all the strength of his faults forgotten, had affected her spirits exceedingly, and thrown her into greater grief for him than she had known on first hearing of his death. Mr Musgrove was, in a lesser degree, affected likewise; and when they reached the cottage, they were evidently in want, first, of being listened to anew on this subject, and afterwards, of all the relief which cheerful companions could give them.

To hear them talking so much of Captain Wentworth, repeating his name so often, puzzling over past years, and at last ascertaining that it \textit{might}, that it probably \textit{would}, turn out to be the very same Captain Wentworth whom they recollected meeting, once or twice, after their coming back from Clifton—a very fine young man—but they could not say whether it was seven or eight years ago, was a new sort of trial to Anne's nerves. She found, however, that it was one to which she must inure herself. Since he actually was expected in the country, she must teach herself to be insensible on such points. And not only did it appear that he was expected, and speedily, but the Musgroves, in their warm gratitude for the kindness he had shewn poor Dick, and very high respect for his character, stamped as it was by poor Dick's having been six months under his care, and mentioning him in strong, though not perfectly well-spelt praise, as »a fine dashing felow, only two perticular about the schoolmaster,« were bent on introducing themselves, and seeking his acquaintance, as soon as they could hear of his arrival.

The resolution of doing so helped to form the comfort of their evening.